\section{Implementation}\label{sec:implementation}
\todo[inline]{Write Implementation transition}

\subsection*{Data Sources}\label{subsec:data-sources}
A candidate link is a link whose line-of-sight availability will be checked in the next step.
The central idea is that while each cell site has its own backhaul requirements, it also acts as a relay for backhaul traffic of its neighbors and their neighbors.
Candidate links, therefore, are always between two cell sites.

We use a filtered version of the French National Frequencies Agency (ANFR) database ``Données sur les installations radioélectriques de plus de 5 watts'' (data on radio installations exceeding 5 W). It is unique because it is an official, country-wide, operator-wide, public source of cell sites including their locations and mounting heights.
Other countries, such as Germany or the United States, do not publish or collect this data, because of national security concerns and different statutory requirements.
Coordinates are given of emitter supports, which are structures on which several antennas by different operators can be mounted.
Around 3/4 are buildings, 1/4 towers and masts, and only 2.4\% of different nature.
While the database names exact coordinates for every support, particularly buildings and wide broadcasting towers are large structures where antennas may be mounted in many locations on the support.
The locations therefore do not provide the exact constraints for where FSO transceivers can be mounted on a support, but only one point.
We further see that when compared to Google Maps and Street View imagery, some database locations appear visibly shifted from the real location of the antenna support.
For all building supports, we run a heuristic improvement by comparing their location to an official public building atlas and snapping it to the closest building.
A quantitative analysis of the error and the heuristic follows in~\cref{subsec:cell-site-locations}.

Simulating realistic lines-of-sight requires a model of obstructions of the world.
We deliberately decide against statistical models, such as the Manhattan Grid Model, in favor of sensor-based models.

OpenStreetMap, as well as the national mapping agency of France (IGN), which is the country of origin of the cell location database, have datasets of building footprints along with a single height value.
This approach has been used in previous works, see~\cref{subsec:related-work}.
Irrespective of their shape, with this method, buildings are limited to be modeled as cuboids, and other obstructions are not captured, such as vegetation or bridges.
Both data sources also lack values for some buildings around Paris.

Commercial options exist, offering increased accuracy based on LiDAR measurements, as well as raw LiDAR point clouds, post-processed formats thereof and multi-sensor measurements.
Availability in France varies.
The evaluated options employ unpredictable, recurring or undisclosed pricing models\todo{should I refer to the tested options here through citations?}.

IGN also offers the public and free LiDAR HD program.
It covers almost the entirety of metropolitan France at $10\ impulses/m^2$ with a total of 4 PB of data\todo{quote?}.

A Digital Surface Model (DSM) is a map raster where each pixel encodes the height of the first and highest LiDAR return, capturing the surface elevation of the covered area.
IGN provides a DSM derived product based on their LiDAR HD program.

We choose this product because DSMs are a standardized, non-proprietary format which may have public data sources in other countries.
It also reduces the data to a more manageable amount by around three orders of magnitude compared to the raw LiDAR point cloud, while maintaining high resolution.

\subsection*{Algorithm}\label{subsec:algorithm}
Conceptually, computing collisions of a line with the obstruction model requires calculating which pixels of the DSM raster it crosses and checking if the surface at any of those pixels is higher than the line.

Line rasterization algorithms from computer graphics, such as Bresenham's and Xiaolin Wu's algorithm, solve a similar problem of painting pixels between two lines~\cite{bresenham, xiaolinwu}.
Designed for another purpose, however, they do not paint exactly all pixels touched by the line, skipping over potential obstacles\todo{show? figure? or good without deeper explanation?}.
Sampling points along the line at a rate of two samples per pixel shows accuracy gains at similar speeds, but still misses pixels.
The Python library rasterio offers a function masking exactly all pixels touched by a line, which in turn calls to GDAL~\cite{gdal}.
The fastest, accurate, CPU-based method tested is ``A Fast Voxel Traversal Algorithm'' by Amanatides and Woo~\cite{amanatides1987fast}.
Hardware accelerators may also be used, but early-exit performance gains will suffer from shared program counters.

\todo{Write \& shorten OptiX talk}
NVIDIA and other graphics processing unit producers offer dedicated ray tracing offloading cores and SDKs. For instance, NVIDIA RTX GPUs can be programmed with the OptiX SDK by building a BVH (bounding volume hierarchy) acceleration structure (AS) in memory and running line-of-sight checks against that almost entirely in hardware.
Our tests show that the main bottleneck becomes the memory transfer of the DSM to the accelerator, the time spent on building the AS and the additional memory occupied by the AS.
The DSM can be seen as a two-dimensional grid of ``pixel pillars'' - each pixel a cuboid stretching from -inf to the surface height at the pixel.
Converting the DSM to triangle primitives without quality loss requires two triangles per visible surface, which are laid out as one rectangle.
Since the top surface of a pixel is always visible, and the bottom pixel is always covered by the sides of this and neighboring pixels and is never visible, and assuming that on average two sides are visible per pixel, this results in six triangles per pixel.
... 3 floats per triangle, 1 after compaction, too many values, memory inflation ...
Using instance acceleration structures (IAS), which allow replicating one AS by shifting and scaling it, may reduce this memory inflation.
NVIDIA supports custom primitives for its ray tracing platform for any objects that cannot be modeled adequately with triangles.
Custom primitives assign an axis-aligned bounding box (AABB) containing the primitive and a custom shader, which will be invoked when a ray intersects the AABB. The ray tracing check is still offloaded and only hands over to the shader for an intersection with the AABB.
Since aforementioned ``pixel pillars'' are AABBs, they can be transformed to custom primitives.
Whenever the intersection shader is called, it already marks an intersection without further computation, and the line can immediately be marked as blocked.
Memory inflation can also be combated by smoothing out the model creating a triangulated RSG (regular square grid).
Each pixel is converted to one triangle that
... conservative conversion (always the maximum) ...
... performance: 6-8 gigarays per second for one 1*1km tile, but lines have to split up over tiles because AS cannot grow much larger ...\
\todo{mention more explicitly that the CPU version generally is line-bound, whereas the RT version is area-bound}

The choice of implementation between CPU, CUDA and OptiX depends primarily on the size of the datasets and region considered, the density of lines per area, implementation complexity and hardware accelerator availability.
For the data and area considered in this paper, the CPU version shows sufficient performance, but derived work with significantly more lines to check will benefit from a hardware ray-traced version.

\subsection*{Adaptability}\label{subsec:adaptability}
\todo[inline]{Write or remove adaptability}